\clearpage
\section{Evaluation}
Looking at \fixme{table}, we can see that the deep-learning model produced significant improvements across all metrics in the 
entire test set. Both TCN Gyro and TCN + MUKF saw around a $38\%$ improvement in mean geodesic error and $40\%$ mean improvement in AYE, when compared to their raw 
gyroscope measurement counterparts. These metrics directly support the fact that the TCN model is able to address drift in tilt and heading
whilst both accelerometer and especially magnetometer 
are disturbed, which are the typical failure modes seen in orientation estimation with IMUs.

Furthermore, comparing TCN dead reckoning to raw gyroscope + MUKF, we can see that the TCN performs better, which highlights that 
the TCN alone does beat uncorrected raw gyroscope measurments even with corrections from an accelerometer and magnetometer. The comparison between
the two produces significant results with improvements in mean geodesic error  $14.6\%$, $31.6\%$ in  mean AYE, and $35.9\%$ in median AYE. This result
is expected as the test set does contain disturbed accelerometer and magnetometer conditions. 

Moreover, median improvements in TCN only of $44\%$ in geodesic error and $53\%$ in AYE, and $39.4\%$, and  $41.9\%$ in TCN + MUKF respectively, show 
a promising result that not only does the TCN produce fewer errors (mean) but also a more consistent output as well. Coupled with the P90 improvements of 
$41.2\%$ in geodesic error, and $44.1\%$ in AYE, for the TCN + MUKF, these results show that the entire distrubution landscape of errors is tighter. 
It shows that not only do we have overall better orientation estimates, but these improvements are more consistent and tail errors are supressed. 


\subsection{Trial 9: Undisturbed, Fast Rotation}
Looking at \fixme{table}, we can see that overall there have been significant improvements in all geodesic error metrics across the two pipelines, with the 
TCN + MUKF pipeline performing the best with an improvement of $37.2\%$ in mean geodesic error and the largest improvements in P90 and standard deviation of $54.8\%$.
This means that the model is able to provide meaningful corrections and that it is robust to tail errors. However, in the TCN Gyro pipeline we can see that
AYE actually gets worse ($-17.7\%$). 


\subsection{Trial 25: Disturbed Tapping}

\subsection{Trial 31: Disturbed Stationary Magnet}

\subsection{Trial 34: Disturbed Magnet attached 3cm}

\subsection{Trial 39: Disturbed, Mixed Conditions}

\subsection{Limitations of the Proposed Solution}

\subsection{Summary}